% !TEX root = main.tex
\clearpage
\section{Code guide: Evaluation and Training Stack}

After cleaning there are now three distinct layers:
\begin{itemize}
    \setlength{\itemsep}{0pt}
  \item tokenize each domain pool;
  \item build the weighted training stream; and
  \item train/evaluate the model.
\end{itemize}

\begin{lstlisting}[style=shadedpython,language={},basicstyle=\scriptsize\ttfamily]
CLEANED DATA POOLS
pools/<domain>/clean.jsonl
        |
        v
tokenize_pool.py
        |
        | uses tokenizer.json
        v
data/domains/<domain>.bin
        +
<domain>.bin.meta.json
        |
        v
build_stream.py
        |
        | applies domain weights
        | holds out validation
        | repeats small pools up to 4x
        v
data/navya0_v1/
|-- train.bin
|-- val.bin
|-- meta.json
        |
        v
train.py
        |
        | reads token batches
        | trains GPT model
        v
checkpoint.pt
        |
        v
run_eval.py
        |
        v
BachattBench scores
\end{lstlisting}

\texttt{build\_stream.py} is the real token-level realization of the mixing
plan. \texttt{mix.py} is more of a planning/manifest stage.

\subsection{\texttt{tokenize\_pool.py}}

Its responsibility is to convert one cleaned domain pool from JSONL text into
a flat stream of BPE token IDs.

Example command:
\begin{lstlisting}[style=shadedpython]
python tokenize_pool.py \
    --pool ../02-data-pipeline/pools/global_english/clean.jsonl \
    --tokenizer ../01-tokenizer/tokenizer-v0.3-64k.json \
    --out data/domains/global_english.bin
\end{lstlisting}

Flow:
\begin{lstlisting}[style=shadedpython,language={},basicstyle=\scriptsize\ttfamily]
global_english/clean.jsonl
        |
        v
read one document
        |
        v
extract ["text"]
        |
        v
tokenizer.encode(text)
        |
        v
token IDs
        |
        v
append <|eos|>
        |
        v
write binary tokens
        |
        v
global_english.bin
\end{lstlisting}

\subsubsection{Why \spectok{<|eos|>} is appended}

This code:
\begin{lstlisting}[style=shadedpython]
buf.extend(tok.encode(text).ids)
buf.append(eos)
\end{lstlisting}
means the model can learn: ``This document ended here.'' Without it, two
unrelated documents could get glued together as though they were continuous
text.

\subsubsection{Why \texttt{tokenize\_pool.py} streams}

This is an important difference from \texttt{prepare\_data.py}.

\texttt{prepare\_data.py} does:
\begin{lstlisting}[style=shadedpython]
ids = []
\end{lstlisting}
and keeps accumulating token IDs in RAM. That is fine for small datasets. But
for 1 billion tokens, keeping everything in a Python list is impractical.

So \texttt{tokenize\_pool.py} buffers only about 1{,}000{,}000 token IDs, then
writes them to disk:
\begin{lstlisting}[style=shadedpython]
if len(buf) >= 1_000_000:
    np.array(buf, dtype=dtype).tofile(out)
\end{lstlisting}

\subsubsection{\texttt{.bin.meta.json}}

For each binary file it also creates metadata:
\begin{lstlisting}[style=shadedpython,language={},basicstyle=\scriptsize\ttfamily]
global_english.bin
global_english.bin.meta.json
\end{lstlisting}
The metadata contains:
\begin{lstlisting}[style=shadedpython]
{
  "pool": "...clean.jsonl",
  "tokenizer": "...tokenizer.json",
  "vocab_size": 64000,
  "dtype": "uint16",
  "docs": 12345,
  "tokens": 123456789
}
\end{lstlisting}
This tells later stages which tokenizer was used, how large the vocabulary
was, what integer dtype was written, how many docs were ingested, and how
many tokens were produced.

\subsection{\texttt{prepare\_data.py}}

This is similar, but simpler and intended for smaller runs.

Example:
\begin{lstlisting}[style=shadedpython]
python prepare_data.py \
    --input "texts/*.txt" \
    --tokenizer tok.json \
    --out data/run1
\end{lstlisting}

It does:
\begin{lstlisting}[style=shadedpython,language={},basicstyle=\scriptsize\ttfamily]
*.txt files
     |
     v
read everything
     |
     v
tokenize
     |
     v
put ALL token IDs in RAM
     |
     v
split last 10% as validation
     |
     v
train.bin
val.bin
meta.json
\end{lstlisting}

The important difference:
\begin{center}
  \renewcommand{\arraystretch}{1.15}
  \small
  \begin{tabular}{@{}p{0.42\textwidth}p{0.48\textwidth}@{}}
    \textbf{\texttt{prepare\_data.py}} &
    \textbf{\texttt{tokenize\_pool.py}} \\
    \hline
    Reads \texttt{.txt} files &
    Reads cleaned \texttt{.jsonl} pools \\
    Accumulates IDs in RAM &
    Streams IDs to disk \\
    Small/research datasets &
    Billion-token-scale pools \\
    Creates final train/val directly &
    Creates one \texttt{.bin} per domain \\
    Simple path &
    Scalable domain-mixing path
  \end{tabular}
\end{center}

So for smoke tests, \texttt{prepare\_data.py} is enough. For real domain
mixing:
\begin{lstlisting}[style=shadedpython,language={},basicstyle=\scriptsize\ttfamily]
tokenize_pool.py
  -> build_stream.py
\end{lstlisting}

\subsection{\texttt{build\_stream.py}}

This is a very important script. Its responsibility is to take the individually
tokenized domain pools and physically construct the final weighted
\texttt{train.bin} and \texttt{val.bin}.

Input:
\begin{lstlisting}[style=shadedpython,language={},basicstyle=\scriptsize\ttfamily]
data/domains/
|-- global_english.bin
|-- global_english.bin.meta.json
|-- indian_english.bin
|-- indian_english.bin.meta.json
|-- indian_languages.bin
|-- ...
\end{lstlisting}

Then:
\begin{lstlisting}[style=shadedpython]
python build_stream.py \
    --domains data/domains \
    --out data/navya0_v1 \
    --total-tokens 1.3e9
\end{lstlisting}

\subsubsection{Tokenizer compatibility}

It first checks tokenizer compatibility. It reads every domain metadata file:
\begin{lstlisting}[style=shadedpython]
vocab = {m["vocab_size"] for m in metas.values()}
dtype = {m["dtype"] for m in metas.values()}
\end{lstlisting}
Then:
\begin{lstlisting}[style=shadedpython]
assert len(vocab) == 1 and len(dtype) == 1, "mixed tokenizers!"
\end{lstlisting}
This is extremely important. Imagine an English bin where token ID 523 means
\texttt{"India"}, and a Hindi bin where token ID 523 means \dev{का}, because
they were generated with different tokenizers. Combining those streams would
be meaningless. All domain bins must use the exact same vocabulary layout and
dtype.

\subsubsection{Validation before mixing}

Validation is separated before mixing. For every domain:
\begin{lstlisting}[style=shadedpython]
n_val = max(1, int(len(arr) * VAL_FRAC))
\end{lstlisting}
where \texttt{VAL\_FRAC = 0.005}, so 0.5\% of each domain is reserved for
validation. Then:
\begin{lstlisting}[style=shadedpython]
train_pool = arr[:-n_val]
val_pool   = arr[-n_val:]
\end{lstlisting}
Example: if global English has 100M tokens, the first 99.5M become the
training pool and the last 0.5M the validation pool. Importantly, validation is
removed before repetition, so validation tokens are not accidentally copied
into training.

\subsubsection{Domain weights}

Weights:
\begin{center}
  \renewcommand{\arraystretch}{1.15}
  \small
  \begin{tabular}{@{}l r@{}}
    \textbf{Domain} & \textbf{Weight} \\
    \hline
    global English & 50\% \\
    Indian English & 17.5\% \\
    Indian languages & 12.5\% \\
    code & 10\% \\
    finance/econ/law & 7.5\% \\
    math/reasoning & 2.5\%
  \end{tabular}
\end{center}
Suppose the total requested is 1.3B tokens. For global English,
$0.50\times 1.3\text{B}=650\text{M}$. For Indian languages,
$0.125\times 1.3\text{B}=162.5\text{M}$. In code:
\begin{lstlisting}[style=shadedpython]
target = w * args.total_tokens
\end{lstlisting}

\subsubsection{Repetition}

Then:
\begin{lstlisting}[style=shadedpython]
epochs = min(
    target / len(train_pool),
    MAX_EPOCHS
)
\end{lstlisting}
This is the same idea as \texttt{mix.py}, but now using real BPE token counts.
Example: an Indian-language train pool of 80M tokens with a desired target of
162.5M gives
\[
  \text{epochs}=162.5/80\approx 2.03,
\]
meaning use the entire pool twice, plus about 3\% of it again.

\subsubsection{Constructing the stream}

This is where \texttt{build\_stream.py} differs fundamentally from
\texttt{mix.py}.
\texttt{mix.py} says ``use 2.03 epochs''; \texttt{build\_stream.py} actually
performs it:
\begin{lstlisting}[style=shadedpython]
full, rem = divmod(got, len(train_pool))
\end{lstlisting}
For example, with \texttt{got = 162.5M} and \texttt{pool = 80M},
\texttt{full = 2} and \texttt{rem = 2.5M}. Then:
\begin{lstlisting}[style=shadedpython]
for _ in range(full):
    train_pool.tofile(train_f)
\end{lstlisting}
writes $80\text{M}+80\text{M}$, and
\begin{lstlisting}[style=shadedpython]
train_pool[:rem].tofile(train_f)
\end{lstlisting}
adds 2.5M. Result: 162.5M tokens actually written into final
\texttt{train.bin}.

\subsubsection{\texttt{mix.py} vs \texttt{build\_stream.py}}

\begin{center}
  \renewcommand{\arraystretch}{1.15}
  \small
  \begin{tabular}{@{}p{0.42\textwidth}p{0.48\textwidth}@{}}
    \textbf{\texttt{mix.py}} &
    \textbf{\texttt{build\_stream.py}} \\
    \hline
    Planning layer &
    Execution layer \\
    Uses supplied pool token counts &
    Reads actual \texttt{.bin} token streams \\
    Calculates epochs &
    Calculates + physically repeats data \\
    Produces snapshot manifest &
    Produces \texttt{train.bin} / \texttt{val.bin} \\
    Does not tokenize &
    Input is already tokenized \\
    Does not physically concatenate &
    Actually writes final streams
  \end{tabular}
\end{center}

Mental model:
\begin{itemize}
    \setlength{\itemsep}{0pt}
  \item \texttt{mix.py} $\rightarrow$ ``This is what we WANT.''
  \item \texttt{build\_stream.py} $\rightarrow$ ``Actually BUILD it.''
\end{itemize}


